\section{System Architecture Overview}
\label{sec:arch}

\subsection{Layered Decomposition}

\EM{} is decomposed into five layers with strictly downward dependencies. No
layer may reference a type defined in a layer above it; this is enforced by
placing each layer in a separate Swift module with explicit public interfaces.

\begin{table}[htbp]
\caption{Module Layering}
\label{tab:layers}
\begin{center}
\renewcommand{\arraystretch}{1.15}
\begin{tabularx}{\columnwidth}{@{}l>{\raggedright\arraybackslash}X@{}}
\toprule
\textbf{Module} & \textbf{Responsibility} \\
\midrule
\texttt{EmulsionApp} & SwiftUI scenes, navigation, app lifecycle. Depends on everything below. \\
\texttt{EmulsionUI} & Reusable views, control surfaces, gesture handling, the Metal-backed preview view. No business logic. \\
\texttt{EmulsionCore} & Domain model: \texttt{Recipe}, \texttt{Stage}, \texttt{Parameter}, \texttt{Asset}, edit history, preset resolution. Pure Swift, no rendering, no I/O. Fully unit-testable on any platform. \\
\texttt{EmulsionRender} & Render graph, resource pooling, Metal pipeline state cache, the pass implementations, LUT baking. Consumes a \texttt{Recipe}, produces textures. \\
\texttt{EmulsionStore} & Persistence: SQLite schema, asset library, cache management, sidecar import/export. \\
\texttt{EmulsionCapture} & AVFoundation session management, ProRAW capture, device capability probing, metadata extraction. \\
\bottomrule
\end{tabularx}
\end{center}
\end{table}

The critical architectural decision is that \texttt{EmulsionCore} contains no
Metal and no Foundation persistence types. The parameter model, the pipeline
topology, and the semantics of every stage are expressed as plain Swift values.
This makes the entire imaging specification testable without a GPU, permits a
reference CPU implementation for validation (Section~\ref{sec:validation}), and
means the macOS companion application in the roadmap shares the domain layer
unchanged.

\subsection{The Development Pipeline}

Figure~\ref{fig:pipeline} shows the full data flow. The pipeline is a directed
acyclic graph, but for v1.0 it is a linear chain with optional bypasses; the
graph machinery exists because the roadmap requires node-based editing and
retrofitting it later would be a rewrite.

\begin{figure*}[htbp]
\centering
\resizebox{\textwidth}{!}{%
\begin{tikzpicture}[
  font=\footnotesize,
  node distance=0pt,
  stg/.style={draw=emLine, rounded corners=2pt, fill=emGray, minimum height=8.5mm,
              text width=20mm, align=center, inner sep=2pt},
  spatial/.style={stg, fill=emAcc!12, draw=emAcc!70},
  dom/.style={font=\scriptsize\itshape, text=emLine},
  ar/.style={-{Latex[length=1.6mm]}, draw=emLine, thick}
]

\node[stg] (s0) at (0,0) {RAW\\decode};
\node[stg, right=5mm of s0] (s1) {Sensor\\normalize};
\node[stg, right=5mm of s1] (s2) {Exposure\\/ WB};
\node[spatial, right=5mm of s2] (s3) {\textbf{Halation}\\(pre-exposure)};
\node[stg, right=5mm of s3] (s4) {Negative\\density};
\node[stg, right=5mm of s4] (s5) {Development\\modifiers};
\node[spatial, right=5mm of s5] (s6) {\textbf{Interlayer}\\inhibition};
\node[spatial, right=5mm of s6] (s7) {\textbf{Grain}\\synthesis};

\foreach \a/\b in {s0/s1, s1/s2, s2/s3, s3/s4, s4/s5, s5/s6, s6/s7}
  \draw[ar] (\a) -- (\b);

\node[stg, below=11mm of s0] (t0) {Print\\exposure};
\node[stg, right=5mm of t0] (t1) {Printer\\lights};
\node[stg, right=5mm of t1] (t2) {Print stock\\density};
\node[stg, right=5mm of t2] (t3) {Color\\science};
\node[spatial, right=5mm of t3] (t4) {\textbf{Optical}\\simulation};
\node[spatial, right=5mm of t4] (t5) {\textbf{Aging}\\/ artifacts};
\node[stg, right=5mm of t5] (t6) {Output\\transform};
\node[stg, right=5mm of t6] (t7) {Encode /\\display};

\foreach \a/\b in {t0/t1, t1/t2, t2/t3, t3/t4, t4/t5, t5/t6, t6/t7}
  \draw[ar] (\a) -- (\b);

\draw[ar] (s7.east) -- ++(0.5,0) |- ($(t0.west)+(-0.8,0)$) -- (t0.west);

\node[dom, above=1.2mm of s0] {DNG};
\node[dom, above=1.2mm of s2] {linear scene};
\node[dom, above=1.2mm of s4] {$\logE \to D$};
\node[dom, above=1.2mm of s7] {negative $D$};
\node[dom, below=1.2mm of t0] {$T = 10^{-D}$};
\node[dom, below=1.2mm of t2] {$\logE' \to D'$};
\node[dom, below=1.2mm of t5] {print $D'$};
\node[dom, below=1.2mm of t7] {Display P3};

\begin{scope}[on background layer]
  \node[fit=(s0)(s7), draw=emBlue!50, dashed, rounded corners=3pt, inner sep=4mm] (negbox) {};
  \node[fit=(t0)(t7), draw=emGreen!50, dashed, rounded corners=3pt, inner sep=4mm] (prbox) {};
\end{scope}
\node[above right=0mm and 0mm of negbox.north west, font=\scriptsize\bfseries, text=emBlue] {NEGATIVE DOMAIN};
\node[below right=0mm and 0mm of prbox.south west, font=\scriptsize\bfseries, text=emGreen] {PRINT DOMAIN};

\end{tikzpicture}}
\caption{The \EM{} development pipeline. Shaded stages are spatial (neighborhood-dependent) or stochastic and cannot be expressed as a lookup table. The pipeline crosses from the negative domain to the print domain by converting density to transmittance, which is the operation a physical enlarger performs.}
\label{fig:pipeline}
\end{figure*}

\subsection{Derivation of Stage Ordering}

The ordering in Figure~\ref{fig:pipeline} is not a preference. Each adjacency
is forced by physics, and the following are the binding constraints:

\begin{enumerate}
\item \textbf{Halation precedes density formation.} Scattered light is light.
It adds to the exposure that forms the latent image, so it must be summed in
the linear exposure domain \emph{before} the characteristic curve is applied.
Applying a glow after tone mapping --- the near-universal shortcut --- produces
halos whose intensity does not correctly saturate, because the shoulder has
already compressed the highlight that should have driven them. This is the
single most visible correctness difference between \EM{} and a filter.
\item \textbf{Interlayer inhibition acts on development, not on the printed
image.} Inhibitors are released during development in proportion to local
development activity, and diffuse before they can suppress neighbors.
Therefore the operator reads the developed density field and modifies it, after
the characteristic curve and before printing.
\item \textbf{Grain lives on the negative.} The dye clouds are in the negative
emulsion. Their density fluctuation is subsequently \emph{printed}, which means
the print stock's contrast amplifies negative grain --- a high-contrast print
stock makes grain more visible, exactly as in practice. Synthesizing grain
after the print transform loses this coupling and produces grain whose
visibility is independent of print contrast, which is wrong and looks wrong.
\item \textbf{Printer lights act on the print exposure, not on the negative
density.} A printer light attenuates the enlarger beam. Because the negative's
transmittance multiplies the beam, and because both are in the exposure domain,
printer lights are a multiplicative offset in linear exposure, equivalently an
additive offset in $\logE$. Section~\ref{sec:printerlights} shows why this
produces the characteristic ``rotation'' of color balance with density that a
simple RGB gain cannot.
\item \textbf{Optical simulation follows the print.} Lens artifacts belong to
the taking lens and therefore ought, strictly, to precede exposure. \EM{}
splits them: \emph{geometric and irradiance} effects (distortion, vignetting,
chromatic aberration) are applied pre-exposure where they physically belong;
\emph{appearance} effects associated with viewing and projection (bloom,
diffusion) are applied post-print. Section~\ref{sec:optical} specifies the
split and the rationale for each.
\item \textbf{Aging is last.} Physical damage to the film or print happens
after everything else and is not attenuated by any subsequent process.
\end{enumerate}

\subsection{Domain Transitions}

The pipeline traverses four distinct numerical domains, and every bug class we
anticipate arises at a transition between them. Table~\ref{tab:domains} is the
normative statement of what quantity is carried at each point.

\begin{table}[htbp]
\caption{Numerical Domains Along the Pipeline}
\label{tab:domains}
\begin{center}
\renewcommand{\arraystretch}{1.15}
\begin{tabularx}{\columnwidth}{@{}l>{\raggedright\arraybackslash}Xll@{}}
\toprule
\textbf{Domain} & \textbf{Quantity} & \textbf{Range} & \textbf{Format} \\
\midrule
Scene linear & Relative exposure $E$, $E=1$ at diffuse white & $[0, 64]$ & RGBA32F \\
Log exposure & $x = \logE$ & $[-4, +2]$ & RGBA32F \\
Negative density & $D$, Status M & $[0.1, 3.5]$ & RGBA16F \\
Print exposure & $E' = T \cdot L$ & $[0,\ 8]$ & RGBA16F \\
Print density & $D'$, Status A & $[0.05, 2.6]$ & RGBA16F \\
Display & Encoded P3 / sRGB & $[0,1]$ & RGBA8 / RGBA16F \\
\bottomrule
\end{tabularx}
\end{center}
\end{table}

Note that the format column permits \texttt{half} precision only from the
density domain onward. In density, values are compressed into roughly
$[0, 3.5]$ with meaningful resolution of $0.001$; \texttt{half} provides
approximately $0.0005$ resolution at $D=2$, which is adequate. In scene linear,
a 16-stop range requires representing $10^{-4}$ and $10^{1}$ with equal relative
precision, and \texttt{half}'s subnormal behavior near $10^{-4}$ introduces
visible shadow quantization. Constraint CR-05 exists for this reason and is
tested by NFR-06.

\subsection{Concurrency Model}

Three execution contexts exist and are strictly separated:

\begin{itemize}
\item \textbf{Main actor.} All SwiftUI state, all user input, all
\texttt{Recipe} mutation. The \texttt{Recipe} is a value type; mutations are
cheap and produce a new immutable snapshot.
\item \textbf{Render actor.} A dedicated serial \texttt{actor} owning the
\texttt{MTLCommandQueue}, the resource heap, and the pipeline state cache.
Receives immutable \texttt{Recipe} snapshots, coalesces them (only the most
recent pending snapshot is rendered --- intermediate ones are dropped), and
enqueues GPU work.
\item \textbf{I/O queue.} A concurrent \texttt{DispatchQueue} for RAW decode,
thumbnail generation, database writes, and export encoding.
\end{itemize}

The coalescing behavior on the render actor is what makes slider dragging feel
correct. During a drag, parameter updates arrive at up to 120\,Hz while the
render takes 8--16\,ms. Queueing every update produces mounting latency; the
render actor therefore holds a single \texttt{pendingRecipe} slot, overwrites
it on each update, and picks up the latest value when the previous frame's
command buffer completes. This is specified in Section~\ref{sec:apparch} with
the full state machine.
